% Chapter Template

\chapter{Conclusions}\singlespacing % Main chapter title

\label{Chapter6} % Change X to a consecutive number; for referencing this chapter elsewhere, use \ref{ChapterX}

\lhead{Chapter 6. \emph{Conclusions}} % Change X to a consecutive number; this is for the header on each page - perhaps a shortened title

%----------------------------------------------------------------------------------------
%	SECTION 1
%----------------------------------------------------------------------------------------

\section{Summary of Research Contributions and Progress}
\label{sec:conclusion_summary}
This report has documented a comprehensive and systematic research program aimed at addressing the critical challenges of enhancing and controlling Large Language Models for low-resource languages. The work has progressed methodically from foundational interpretability and advanced theoretical explorations to direct, practical application, yielding significant insights that have actively shaped the research trajectory.

The research journey commenced with a foundational investigation into the internal mechanisms of multilingual models. This work, presented at the NAACL 2024 Workshop \citep{mondal2025language}, rigorously tested the hypothesis that language-specific neurons could be leveraged for cross-lingual transfer. The key finding—that direct interventions on these neurons are insufficient for improving downstream task performance due to their polysemantic nature—provided a critical, empirical justification for pivoting away from surgical internal edits and towards more holistic, output-level control mechanisms.

To build the necessary expertise for such control, a second foundational project was undertaken, culminating in a full paper submission to NeurIPS 2025 \citep{mondal2025fairpo}. This research, \textit{FairPO}, involved designing a novel, sophisticated framework for fair multi-label classification. It required moving beyond the standard application of preference optimization to adapt its core principles to a new domain, integrate them with group-robust optimization (GRPO), and handle dynamically generated preference sets. This project provided an indispensable, deep, and practical mastery of advanced model alignment frameworks.

The current research phase applies this expertise to the task of low-resource Neural Machine Translation. Preliminary experiments were conducted to fine-tune Llama 3.1 for English-to-Bengali translation using Direct Preference Optimization (DPO). The initial results were unexpected, showing a significant degradation in performance compared to both a one-shot and a standard fine-tuning baseline. This negative result is, in itself, a crucial scientific finding, highlighting the non-trivial complexities and potential instabilities of applying preference optimization in this specific context. It has since prompted a necessary and rigorous debugging phase, which is currently underway.

\section{Overall Thesis Contribution and Future Outlook}
\label{sec:conclusion_outlook}
The primary contribution of the work to date is a multi-faceted investigation into model control for low-resource NLP. This includes a peer-reviewed contribution that challenges a prevailing hypothesis in interpretability, a conference-submitted paper that showcases mastery of advanced optimization, and a critical empirical finding that informs the practical application of DPO.

The future work is now well-defined and strategically planned. The immediate priority, is to systematically diagnose the issues encountered in the NMT experiments to produce a conclusive result. Following this consolidation phase, the research will advance to a novel and ambitious objective: the development of a \textit{unified, agentic framework for low-resource QA and translation}.

This proposed framework represents a significant conceptual step forward, aiming to synthesize the strengths of several state-of-the-art paradigms. It will use the agentic, multi-step reasoning pipeline of \textit{KGAREVION} \citep{su2025kgarevion} as its blueprint, adapt it for unstructured text retrieval, and instill robustness against noisy context by incorporating the training principles of \textit{RAFT} \citep{zhang2024raft}. Critically, the final generation module of this agent will be enhanced using the preference optimization techniques from Phase 1 of this research, creating a powerful synergy between factually grounded reasoning and high-quality linguistic generation.




